% Generated draft: requires mathematical and source review.
\documentclass[11pt]{article}
\usepackage[margin=25mm]{geometry}
\usepackage{fontspec}
\setmainfont{Latin Modern Roman}
\usepackage{amsmath,amssymb,mathtools,unicode-math}
\setmathfont{Latin Modern Math}
\usepackage{enumitem,needspace,xurl,hyperref}
\setlength{\parindent}{0pt}
\setlength{\parskip}{5pt}
\setlength{\emergencystretch}{3em}
\newcommand{\N}{\mathbb{N}}
\newcommand{\Z}{\mathbb{Z}}
\newcommand{\Q}{\mathbb{Q}}
\newcommand{\R}{\mathbb{R}}
\newcommand{\C}{\mathbb{C}}
\newcommand{\F}{\mathbb{F}}
\newcommand{\A}{\mathbb{A}}
\newcommand{\PP}{\mathbb{P}}
\newcommand{\E}{\mathbb{E}}
\newcommand{\eps}{\varepsilon}
\newcommand{\dd}{\,\mathrm{d}}
\begin{document}
\section*{UnsolvedMath ID: 801167}
\textbf{Source problem code:} catalogue entry
\subsection{Definitions and mathematical statement}
Let \(n\geq 1\) be an integer and let \(H=L^2(\R^n)\), with Lebesgue measure. A countable family \(\mathcal W\subset H\) is a Parseval frame if
\[
\|f\|_2^2=\sum_{w\in\mathcal W}|\langle f,w\rangle|^2
\qquad\text{for every }f\in H.
\]
A Parseval wavelet is a function \(\psi\in H\) whose associated affine wavelet family is a Parseval frame. Write \(V(\psi)\) for the closed linear span of the negative dilation levels of that family, as intended by the source. The precise affine family and dilation convention require clarification below.

The supplied question is whether every such Parseval wavelet satisfies
\[
\bigcap_{j\in\Z}D^j\bigl(V(\psi)\bigr)=\{0\},
\qquad V_j=D^j\bigl(V(\psi)\bigr)\quad(j\in\Z),
\]
where \(D\) is the intended invertible dilation operator and \(D^j(V)=\{D^jv:v\in V\}\).

\textbf{Unresolved operator specification.} The supplied record explicitly writes
\[
(Df)(x)=\sqrt{2}\,f(x).
\]
This formula contains no dilation of the argument and is not a unitary operator on \(L^2(\R^n)\). The record does not specify a dilation matrix or a replacement argument for \(f\), nor does it give an explicit formula for the affine family or its negative dilation levels. Consequently, the displayed intersection question preserves the intended target, but a fully specified wavelet formulation requires verification of these conventions against the cited source. No corrected operator is selected here.

The supplied background attributes to Baggett the assertion that \(\{V_j:j\in\Z\}\) satisfies the generalized multiresolution analysis properties labelled (M1), (M2), (M4), and (M5'), with (M3) possibly failing. These labels are not defined in the supplied record; this attribution is retained as source context rather than an independently verified theorem.
\subsection{Short English statement}
The question asks whether the intersection of all integer dilation images of the negative-dilate space of every Parseval wavelet is trivial. The supplied dilation formula and wavelet conventions need clarification before this becomes a fully specified mathematical statement. The registry labels the problem open; that status is supplied source data.
\subsection{Sources}
\begin{itemize}
\item Local UnsolvedMath registry, permanent ID 801167, source code catalogue entry, title \emph{Open Problems in Wavelet Theory --- Problem 6}; collection: AMR Open Problem Lists. The record attributes the problem to Bownik and Rzeszotnik, with proposed year 2019, and labels it open.
\item Bownik and Rzeszotnik, \emph{Open Problems in Wavelet Theory} (2019), Problem 6. Supplied source URL: \url{https://arxiv.org/abs/1905.08230}. The provenance records access on 2026-10-16 and extraction by \texttt{arxiv\_latex}; these are supplied provenance statements, not a claim of independent access or verification.
\item The supplied background attributes the stated generalized multiresolution analysis properties to Baggett, but supplies no separate bibliographic reference for that assertion.
\end{itemize}
\end{document}
